%% sections/discussion.tex — Discussion
%% Target: ≥400 words, covering mechanism, contrast with chaos, open questions
\label{sec:discussion}

\subsection{Physical mechanism}
\label{sec:mechanism}

The results presented above point to a coherent physical picture.
At $t = 0$, the product initial state encodes all the information about Z2 symmetry
breaking directly in its local structure: the EA $\DS_A(0;\theta) = \hbin((1+\cos^{\NA}\!\theta)/2)$
is determined entirely by the initial polarisation angle and the subsystem size.
The subsequent quench dynamics drives quasiparticle excitations across the chain,
spreading correlations and washing out the memory of the initial state.
The \emph{rate} at which this spreading restores Z2 symmetry in the subsystem
depends on the distribution of quasiparticle modes that are activated by the quench —
a distribution that is, in turn, sensitive to the initial state through the overlap
between $|\psi_\theta\rangle$ and the eigenstates of the post-quench Hamiltonian.

The key insight revealed by our systematic study is that a larger initial EA imbalance
$|\Delta\Delta\mathcal{S}|$ produces a \emph{later} QME crossing, not an earlier one
(Spearman $\rho = +0.90$).
This is physically reasonable: a larger head start in EA requires more evolution for
the more-broken state to "close the gap" and overtake the less-broken companion.
The crossing time $\tM$ therefore encodes initial-state information through the
imbalance in initial conditions, analogously to how the classical Mpemba effect
encodes the initial temperature imbalance.
However, the specific mechanism by which quasiparticle dynamics translates an initial
EA imbalance into a crossing time — and in particular whether this can be captured by
a perturbative argument at early times — remains an open analytical question.
An early-time expansion of $\DS_A(t)$ around $t=0$ for the non-Gaussian product states
studied here, analogous to the quasi-particle picture used for Gaussian states
in Ref.~\cite{Rylands2024}, would provide a theoretical grounding for the
empirical $\rho=+0.90$ correlation.

\subsection{Contrast with chaotic dynamics}
\label{sec:chaos}

The present results complete a picture suggested by general arguments about
chaotic dynamics.
In random unitary circuits, Z2 QME is expected to be absent: rapid scrambling
makes all initial states with the same $\langle Q_A\rangle$ relax the entanglement
asymmetry at the same rate, regardless of the fine-grained structure of the initial
state, because quasi-local conserved structures that distinguish different initial
states are destroyed at short times.

Our results show the opposite behaviour in integrable free-fermion chains: Z2 QME
is \emph{universal}, occurring for every parameter combination tested.
The integrability of the TFIM (and the XY chain family) ensures the existence of
quasi-local conserved quantities whose overlap with the initial state persists
throughout the dynamics.
Different initial polarisation angles $\theta$ excite different quasi-particle
mode distributions, and the resulting differential relaxation rates produce the
Mpemba crossing.
The contrast (chaotic $\to$ Z2 QME absent; integrable $\to$ Z2 QME universal)
mirrors the known contrast for the classical Mpemba effect: the effect is associated
with specific non-equilibrium initial conditions that are ``remembered'' by the dynamics.
We conjecture that this dichotomy extends beyond the free-fermion case and constitutes
a general feature of integrable versus chaotic symmetry-restoring dynamics.

\subsection{Role of dynamical quantum phase transitions}
\label{sec:dqpt_role}

Dynamical quantum phase transitions are among the most prominent non-equilibrium
phenomena in isolated quantum systems~\cite{Heyl2018}, and it was natural to hypothesise
that they might organise the timing of Mpemba crossings.
Our statistical analysis provides a nuanced answer to this question.

DQPTs do imprint weakly on the EA dynamics: the distribution of fractional
crossing positions $\phi$ is non-uniform (KS $p \approx 6 \times 10^{-25}$),
with crossings tending to occur in the first half of each DQPT period
($\bar\phi = 0.377$, peak at $\phi \approx 0.15$).
This suggests that the DQPT periodicity modulates the \emph{envelope} of
$\DS_A(t)$ oscillations — as the Loschmidt echo passes through successive
minima, the EA curves of different initial states are differentially affected,
preferentially enabling crossings in the rising phase of each Loschmidt oscillation
rather than near the cusp.
However, DQPTs do not \emph{trigger} crossings: there is no excess concentration
near $\phi = 0.5$ (binary test $p = 0.27$), and QME occurs for all 23 post-quench
fields tested including eight FM-phase fields where DQPTs are entirely absent.

This finding clarifies which observables are and are not organised by DQPTs.
While DQPTs manifest directly in the Loschmidt return rate~\cite{Heyl2013} and
can influence two-point order-parameter correlators~\cite{Zunkovic2018},
the EA crossing time $\tM$ is primarily a function of initial state geometry,
placing it in a distinct universality class from DQPT-sensitive observables.

\subsection{Open questions and outlook}
\label{sec:outlook}

Our work opens several directions for future investigation.

\emph{Analytical theory.}
A microscopic theory of $\tM$ for the non-Gaussian product states studied here
would require extending the quasiparticle picture of EA dynamics~\cite{Rylands2024}
to states that break fermion parity.
The parity-sector mixing identified in Sec.~\ref{sec:quench} — the initial state
being a coherent superposition of the even and odd sectors — may enable an
analytical treatment via the Ramond/Neveu-Schwarz sector decomposition.
Such a theory could explain both the $\rho = +0.90$ correlation and the
$\phi \approx 0.15$ peak of the fractional-position distribution.

\emph{Beyond free fermions.}
The free-fermion structure of the TFIM and XY chain ensures exact solvability
but restricts the scope of the results.
Extending the analysis to interacting integrable systems (XXZ chain, Hubbard model)
and to symmetry classes beyond Z2 (e.g.\ $\mathbb{Z}_N$ or $\mathrm{U}(1)$) would
determine whether the Z2 QME universality observed here is specific to free-fermion
dynamics or reflects a broader principle.

\emph{Experimental realisations.}
Transverse-field Ising chains are directly realised in superconducting qubit
arrays~\cite{Bloch2008} and trapped-ion quantum simulators, with the ability to
prepare product spin states of the form~\eqref{eq:psi0} and measure subsystem
entanglement via quantum-state tomography.
Our exact initial formula $\DS_A(0;\theta) = \hbin((1+\cos^{\NA}\!\theta)/2)$
provides an unambiguous prediction for the pre-quench asymmetry, and the
convergence of $\tM$ to a finite thermodynamic-limit value at $N \sim 14$
makes the effect accessible at currently achievable system sizes.

\emph{Mpemba effect engineering.}
The strong $\rho = +0.90$ correlation between $|\Delta\Delta\mathcal{S}|$ and $\tM$
suggests a practical handle: by tuning the initial polarisation angle, one can
engineer the Z2 restoration time over nearly four orders of magnitude without changing
any post-quench parameter.
This tunability may find applications in quantum state-preparation protocols where
the speed of symmetry restoration is a resource.
